% !TEX TS-program = lualatex

%%%
% src::
%     url = https://github.com/pfradin/luadraw/discussions/194#discussioncomment-15836574
%%%

\documentclass[border = 3pt]{standalone}

%%%
% Des couleurs nommées faciles d'emploi via le package \xcolor.
%%%
\usepackage[svgnames]{xcolor}

%%%
% La bibliothèque \luadraw allie une facilité d’utilisation
% à un rendu particulièrement soigné.
%%%
\usepackage[3d]{luadraw}


%%%
% La bibliothèque ''shadings'' proposée par \tikz, déjà chargé
% par \ludadraw, est utilisée par ''luadraw_shadedforms'', une
% extension proposée par \ludadraw.
%%%
\usetikzlibrary{shadings}

\begin{document}

\begin{luadraw*}{name = vector-field-two-grids}
local palext = require 'luadraw_palettes'

require 'luadraw_shadedforms'
require 'luadraw_fields'
require 'luadraw_colors'

local ld = luadraw

local abs = ld.cpx.abs
local Z   = ld.cpx.Z

------
-- Réglages liés au tracé.
------
local xmin, xmax = -3, 3
local ymin, ymax = -3, 3

------------
-- Champs défini par `y' = f(x, y)` où `f(x, y) = x^2 + y^2 - 1`
-- avec la ¨cond initiale `y(0) = 1/2`.
------
local f = function(x, y)
  return x^2 + y^2 - 1
end

local field = function(x, y)
  return {1, f(x, y)}
end

local A_0 = Z(0, 1/2)

------
-- ¨Def du module du champs.
------
local f_modulus = function(x, y)
  local A = field(x, y)

  return abs(Z(A[1],A[2]))
end

------
-- ¨Def de la zone graphique.
------
local graphview = ld.graph:new{
  window = {
    xmin - 0.5, xmax + 1.5,
    ymin - 0.5, ymax + 0.5
  },
  size   = {10, 10},
  bbox   = false
}

------
-- Le fond coloré en fonction de `abs(1 + i(x^2 + y^2 - 1))`.
------
graphview:Dshadedrectangle(
  xmin, xmax, ymin, ymax,
  ld.mixpalette(palext.Jet, 50, ld.White),
  {
    values = f_modulus,
    bar    = "right",
    baroptions
  }
)

------
-- Tracé du champs.
------
graphview:Dvectorfield(
  field,
  {
    view         = {xmin, xmax, ymin, ymax},
    draw_options = "-stealth"
  }
)

------
-- Ajout d'un cadre gradué.
------
graphview:Dgradbox(
  {
    Z(xmin, ymin),
    Z(xmax, ymax),
    1, 1
  },
  {
    title = "Vector field $f(x,y) = (1, x^2 + y^2 - 1)$"
  }
)

------
-- ¨luadraw va résoudre et tracer la ¨sol de l'¨eq_diff.
------
graphview:Dodesolve(
  f,
  A_0.re, A_0.im,
  {
    t            = {xmin, xmax},
    clip         = {-2.5, 2.25, ymin, ymax},
    nbdots       = 150,
    draw_options = "Navy, line width = 1pt"
  }
)

------
-- Ajout du point initial.
------
graphview:Dlabeldot(
  "{\\boldmath $A_0$}", A_0,
  {
    pos          = "S",
    node_options = "Navy"
  }
)

------
-- Montrons le résultat de notre oeuvre.
------
graphview:Show()
\end{luadraw*}

\end{document}
